\newpage

\appendix
\raggedbottom
\section*{Appendix}


\section{Implementation Details}
\label{sec:imple}

\subsection{Dataset Statistics}

\begin{table}[H]
\centering
\footnotesize

\resizebox{\columnwidth}{!}{%
\setlength{\tabcolsep}{3pt}
\begin{tabular}{lrrrrl}
\toprule
Dataset & Total & Train & Val & Test & Freq. \\
\midrule
ETTh1/h2 & 17,420 & 12,194 & 1,742 & 3,484 & Hourly \\
ETTm1/m2  & 69,680 & 48,776 & 6,968 & 11,042 & 15-min \\ 
Weather & 52,696 & 36,887 & 5,269 & 10,540 & 10-min \\
ECL & 26,304 & 18,412 & 2,631 & 5,261 & Hourly \\
Exchange & 7,588 & 5,311 & 759 & 1,518 & Daily \\
\bottomrule
\end{tabular}
}
\caption{Dataset statistics and split configuration.}
\label{tab:datasets}
\end{table}

\subsection{Hyperparameters}
\begin{table}[H]
\centering
\resizebox{\columnwidth}{!}{%
\begin{tabular}{l|ccc}
\toprule
Dataset & LR & Epochs  & STL (Seas., Trend) \\
\midrule
ETTh1 & 1e-4 & 40 & 24, 48 \\
ETTh2 & 1e-4 & 40 & 24, 48 \\
ETTm1 & 1e-4 & 40 & 24, 48 \\
ETTm2 & 1e-4 & 40 & 24, 48 \\
ECL & 5e-6 & 40 & 12, 24 \\
Exchange & 1e-4 & 40 & 7, 24 \\
Weather & 3e-5 & 40 & 7, 48 \\
\bottomrule
\end{tabular}%
}
\vspace{-2mm}
{\footnotesize $^\dagger$Weight decay=1e-7, LLM temp=0.3, top-$p$=0.8 for all.}
\caption{Dataset-specific hyperparameters.$^\dagger$}
\label{tab:hyperparams}
\end{table}


Table~\ref{tab:hyperparams} summarizes dataset-specific hyperparameters for decoder training, STL decomposition, and LLM sampling. The residual decoder is a 2-layer MLP with hidden dimension $L \times F$, scaling capacity with input complexity.
The text projection dimension for the Irregular Agent embedding is set to 32.
All experiments were executed on a system with an NVIDIA RTX 3080 GPU, Intel Xeon E5-2686 v4 CPU (2.30 GHz), and 128GB RAM. Code is available at \url{https://github.com/mfriendly/CTRL}.

\subsection{STL Context for Control Signal Generation}
\label{app:stl_context}
For each sample, agents receive historical (input sequence) and prediction (backbone output) statistics:
\begin{equation}
    \mathcal{C} = \{(\mathcal{C}_\text{hist}, \mathcal{C}_\text{pred}) \mid \mu, \sigma, \text{slope}_T, \text{amp}_S, d_S\}
\end{equation}
where $\mu$ and $\sigma$ denote mean and standard deviation, $\text{slope}_T$ is the trend slope via linear regression, $\text{amp}_S$ is the seasonal amplitude (max $-$ min), and $d_S$ is seasonal dominance (ratio of seasonal variance to total variance). This dual-view enables diagnosis by comparing expected versus predicted temporal characteristics.

\subsection{Shift Detection}
\label{app:shift_detection}

We quantify distribution shift separately for each STL component using z-score normalization against held-out training statistics.

Trend shift captures slope changes:
\begin{equation}
s_\text{trend} = \frac{|\text{slope}_\text{test} - \bar{\text{slope}}_\text{val}|}{\sigma_{\text{slope,val}} + \epsilon}
\end{equation}

Seasonal shift measures dominance changes:
\begin{equation}
s_\text{seasonal} = \frac{|d_{S,\text{test}} - \bar{d}_{S,\text{val}}|}{\sigma_{d_S,\text{val}} + \epsilon}
\end{equation}
where seasonal dominance \textit{$d_S = \text{Var}(\text{seasonal}) / \text{Var}(\text{total})$}.

Irregular shift tracks noise level:
\begin{equation}
s_\text{irregular} = \frac{|\sigma_\text{test} - \bar{\sigma}_\text{val}|}{\sigma_{\sigma,\text{val}} + \epsilon}
\end{equation}

Adaptation activates when $s > \tau$ (default $\tau=2.0$), corresponding to test statistics exceeding two standard deviations from the held-out training distribution. The component-wise formulation allows targeted adjustment: a trend reversal need not affect seasonal correction.


\subsection{Hyperparameter Sensitivity}

\begin{table}[H]

\centering
\small
\begin{tabular}{lc}
\toprule
\textbf{Configuration} & $\Delta$\textbf{MSE (\%)} \\
\midrule
\multicolumn{2}{l}{\textit{LLM Sampling (Llama 3.3 70B)}} \\
\quad temp=0.2, top-p=0.8 & +0.22 \\
\quad \textbf{temp=0.3, top-p=0.8} & \textbf{--} \\
\quad temp=0.5, top-p=0.8 & +0.26 \\
\midrule
\multicolumn{2}{l}{\textit{Adaptation Threshold ($\tau$)}} \\
\quad $\tau$=OFF (no adaptation) & +1.03 \\
\quad $\tau$=1.5 & +0.05 \\
\quad $\boldsymbol{\tau}$\textbf{=2.0} & \textbf{--} \\
\quad $\tau$=2.5 & +0.01 \\
\bottomrule
\end{tabular}

\caption{Hyperparameter sensitivity on ETTh2 (384$\to$96). Bold indicates selected configuration. $\Delta$ shows degradation relative to selected setting.}
\label{tab:hparam_sensitivity}
\end{table}

Table~\ref{tab:hparam_sensitivity} evaluates CTRL's sensitivity to key hyperparameters on ETTh2. LLM sampling parameters (temperature, top-p) have minimal impact across the tested range, indicating that control signal generation is stable across reasonable sampling configurations. The adaptation threshold $\tau$ shows an optimum at 1.5--2.5, corresponding to the z-score formulation where adaptation triggers for test statistics exceeding two standard deviations from held-out training. Disabling adaptation entirely degrades performance by 1.03\%. 


\section{Efficiency Analysis}
\label{app:efficiency}

\subsection{LLM Call Analysis.}
Test-time adaptation monitors distribution shifts at fixed intervals rather than per-sample, with total calls equal to $3 + 3K$ in the worst case where all $K$ checks trigger adaptation (Table~\ref{tab:llm_calls}). Per-sample policies yield only $<$1\% MSE improvement at 200$\times$ cost compared to shared policies, and varying the calibration interval from 5 to 100 batches changes MSE by under 0.15\% across all six datasets (identical on ETTh1). This confirms that Phase 1's initial control signal generation captures the primary distribution gap, while Phase 3 provides marginal refinement, allowing coarse-grained monitoring (default: 50 batches) to reduce LLM calls by 10--50$\times$ without accuracy loss.




\subsection{Model Complexity Breakdown.}
Table~\ref{tab:param_breakdown} provides detailed parameter counts.
The frozen GPT-2 encoder (124M parameters) adds no training cost; 
only the lightweight MLP projector (100K) and residual decoder (215K) are trained.

\begin{table}[H]
\centering

\small
\begin{tabular}{p{4.5cm}p{1.2cm}}
\toprule
\textbf{Component} & \textbf{Params} \\
\midrule
DLinear backbone (frozen) & 74K \\
GPT-2 encoder (frozen) & 124M \\
MLP projector (trained) & 100K \\
Residual decoder wo proj. (trained) & 215K \\
\midrule
\textbf{Total trainable} & \textbf{$\sim$400K} \\
\bottomrule
\end{tabular}
\caption{CTRL parameter breakdown (ETTh2, DLinear backbone).}
\label{tab:param_breakdown}
\end{table}








\section{Case Studies}
\label{app:case_studies}

\subsection{Quantitative: Ablation Analysis}
\label{app:ablation_analysis}

\begin{figure*}[h!]
    \centering
    \includegraphics[width=\linewidth]{fig/grid.png}
\caption{Three-way ablation on ETTh2. Scatter plots (top): zero control signal clusters around the diagonal with near-random win rate (49\%), indicating the decoder alone learns minimal correction. Adding LLM control signals shifts points below the diagonal, but gains concentrate in low-to-mid MSE regions. The full system extends improvements to harder samples where distribution shift invalidates cached control signals. Histograms (bottom): test-time adaptation compresses negative-skewed tails by detecting when cached control signals fail and triggering re-reasoning.}
    \label{fig:3way_ablation}
\end{figure*}

We conduct a three-way ablation to isolate contributions of each CTRL component: (1) Zero Control Signal (decoder only, null control signal), (2) w/o Test-Time Adaptation (LLM-generated control signal only), and (3) CTRL Full (control signal generation + test-time adaptation). Results are evaluated on ETTh2 (384$\to$96) with DLinear backbone.

\begin{table}[h!]
\centering
\footnotesize
\begin{tabular}{lrrr}
\toprule
& \textbf{Zero} & \textbf{w/o TTA} & \textbf{Full} \\
\midrule
Mean Imp. & +4.3\% & +12.6\% & +13.3\% \\
Win Rate & 49\% & 76\% & 82\% \\
\bottomrule
\end{tabular}
\caption{Three-way ablation on ETTh2.}
\label{tab:3way_ablation}
\end{table}

\begin{table}[h!]
\centering
\small
\footnotesize
\begin{tabular}{lc}
\toprule
\textbf{Comparison} & \textbf{$\Delta$ Improv.} \\
\midrule
w/o Adaptation vs Zero & +8.3\% \\
CTRL (Full) vs w/o Adaptation & +0.7\% \\
CTRL (Full) vs Zero & +9.0\% \\
\bottomrule
\end{tabular}
\caption{Relative contribution of each component.}
\label{tab:3way_contributions}
\end{table}

Table~\ref{tab:3way_ablation} quantifies overall performance, but the aggregate metrics obscure where each component provides value. The near-random win rate (49\%) under zero control signal confirms that without LLM guidance, the residual decoder lacks sufficient inductive bias to systematically correct backbone errors; the decoder architecture alone cannot discover correction strategies from data.

The 27\% win rate jump when adding LLM control signals (49\%$\to$76\%) reflects LLM reasoning's capacity to diagnose backbone failure modes from held-out training statistics and encode correction strategies into compact control signals. However, Figure~\ref{fig:3way_ablation} reveals these gains concentrate disproportionately on easier samples (Q1-Q2), where backbone errors follow patterns visible in held-out training data. On harder samples (Q3-Q4), improvements without adaptation diminish because the control signal assumes distributional stationarity that does not hold.

Test-time adaptation provides the critical bridge to hard samples. While its aggregate contribution appears modest (+0.7\% mean improvement, +6\% win rate), this masks its selective impact: on Q4 samples where the system without adaptation shows only +1.4\% improvement, CTRL (Full) achieves +5.7\%. The mechanism is straightforward: by monitoring STL statistics at test time, CTRL detects when input distributions diverge from held-out training and triggers re-reasoning rather than applying stale control signals. This explains the histogram compression in Figure~\ref{fig:3way_ablation} (bottom right), where the negative tail from hard-sample failures shrinks under full adaptation.

The decomposition in Table~\ref{tab:3way_contributions} thus understates test-time adaptation's role. LLM control signal generation accounts for the majority of aggregate improvement (+8.3\% of +9.0\%), but test-time adaptation determines whether CTRL degrades gracefully or catastrophically under distribution shift. For deployment scenarios where test distributions are unpredictable, this robustness property may matter more than average-case gains.

\subsection{Qualitative: Agent Reasoning Traces}

We examine two contrasting cases from ETTh2 (384$\rightarrow$96) to illustrate how test-time adaptation responds to different shift patterns. We select samples exhibiting distinct shift patterns: Sample 0 (highest trend shift score in batch 0) and Sample 640 (highest seasonal shift score in batch 1).

\subsubsection{Case 1: Trend Reversal}

The trend slope reverses from $+3.05 \times 10^{-5}$ (held-out training) to $-1.57 \times 10^{-5}$ (test), yielding a shift score of 1.51 (151\%). Seasonal dominance increases moderately from 0.395 to 0.486 (23\% shift), while noise level rises slightly from $\sigma=1.07$ to $\sigma=1.19$ (11\% shift). The held-out training control signal assumed a rising trend; applying it to declining test data would systematically overpredict.

\begin{figure}[h!]
\centering
\fbox{\parbox{0.95\columnwidth}{
\textbf{Trend Agent} \\
\textit{Diagnosis:} The validation and test trends have a significant shift in direction with a ratio of 151.33\%, indicating a reversal in trend. \\
\textit{Reasoning:} Reduced scale multiplier to dampen the impact of extreme predictions, negative bias shift to adjust for the reversal in trend direction, increased gate multiplier to allow more flexibility in response to new data, and a confidence penalty to reflect reduced certainty due to the significant trend change.
}}
\caption{Trend agent reasoning for Case 1 (trend reversal).}
\label{fig:case1_trend}
\end{figure}

\begin{figure}[h!]
\centering
\fbox{\parbox{0.95\columnwidth}{
\textbf{Seasonal Agent} \\
\textit{Diagnosis:} The test data shows a stronger seasonal pattern compared to the validation data with a significant shift ratio. \\
\textit{Reasoning:} Adjustments are made to amplify the seasonal scale and gate multipliers to better capture the increased strength of the seasonal pattern in test data, while a slight bias shift is added to align predictions with observed trends. Confidence penalty reduces overconfidence due to model adjustments.
}}
\caption{Seasonal agent reasoning for Case 1.}
\label{fig:case1_seasonal}
\end{figure}

Table~\ref{tab:case1_policy} shows the resulting numerical control signal changes. Trend parameters shift dramatically (scale halved from 0.90 to 0.50, confidence reduced from 0.60 to 0.21), while seasonal scale increases from 0.95 to 1.50 to amplify the strengthening pattern. Gates increase across all components (0.40$\rightarrow$0.69 for trend), reflecting higher uncertainty where the LLM delegates more correction authority to the learned decoder.

\begin{table}[h!]
\centering
\small


\begin{tabular}{lcccc}
\toprule
& Scale & Bias & Gate & Conf. \\
\midrule
\multicolumn{5}{l}{\textit{Before Adaptation (Stage 1)}} \\
Trend & 0.90 & $+$0.10 & 0.40 & 0.60 \\
Seasonal & 0.95 & $+$0.10 & 0.40 & 0.60 \\
Irregular & 0.95 & $+$0.10 & 0.50 & 0.60 \\
\midrule
\multicolumn{5}{l}{\textit{After Adaptation (Stage 3)}} \\
Trend & 0.50 & $-$0.50 & 0.69 & 0.21 \\
Seasonal & 1.50 & $+$0.55 & 0.66 & 0.41 \\
Irregular & 1.26 & $-$0.20 & 0.67 & 0.31 \\
\bottomrule
\end{tabular}
\caption{Control signal values before and after adaptation (Case 1: Trend Reversal).}
\label{tab:case1_policy}
\end{table}


\subsubsection{Case 2: Seasonal Amplification (Sample 640)}

Here the pattern inverts: trend shift is minimal (score 0.02, slope unchanged at $\approx 3.0 \times 10^{-5}$), but seasonal dominance more than doubles from 0.395 to 0.849 (115\% shift). The test window falls in a period where periodic fluctuations overpower the stable trend.

\begin{figure}[h!]
\centering
\fbox{\parbox{0.95\columnwidth}{
\textbf{Trend Agent} \\
\textit{Diagnosis:} A small positive shift in the trend slope has been observed with a minimal change ratio of 2.24\%, indicating a slight increase in the rate of change without any direction reversal. \\
\textit{Reasoning:} Adjustments are made to slightly increase the scale and gate multipliers to account for the positive shift in trend slope, while a small bias shift is applied to align predictions more closely with observed patterns.
}}
\caption{Trend agent reasoning for Case 2 (seasonal amplification).}
\label{fig:case2_trend}
\end{figure}

\begin{figure}[h!]
\centering
\fbox{\parbox{0.95\columnwidth}{
\textbf{Seasonal Agent} \\
\textit{Diagnosis:} The test data shows a significantly stronger seasonal pattern compared to the validation data. \\
\textit{Reasoning:} Increasing the scale multiplier and gate multiplier to account for the stronger seasonal pattern in test data, adding a bias shift to further adjust predictions upwards, and applying a confidence penalty to reflect increased uncertainty in seasonal predictions due to the significant pattern change.
}}
\caption{Seasonal agent reasoning for Case 2.}
\label{fig:case2_seasonal}
\end{figure}

Table~\ref{tab:case2_policy} shows the contrasting adaptation. Unlike Case 1, trend parameters change minimally (scale 0.90$\rightarrow$0.58, moderate adjustment). The seasonal agent, detecting strong amplification, maximizes scale to 1.50 and increases bias to $+$0.60 to match the new dominance level. This demonstrates context-sensitivity: the same adaptation mechanism produces different adjustments depending on which component exhibits shift.

\begin{table}[h!]
\centering
\small

\begin{tabular}{lcccc}
\toprule
& Scale & Bias & Gate & Conf. \\
\midrule
\multicolumn{5}{l}{\textit{Before Adaptation (Stage 1)}} \\
Trend & 0.90 & $+$0.10 & 0.40 & 0.60 \\
Seasonal & 0.95 & $+$0.10 & 0.40 & 0.60 \\
Irregular & 0.95 & $+$0.10 & 0.50 & 0.60 \\
\midrule
\multicolumn{5}{l}{\textit{After Adaptation (Stage 3)}} \\
Trend & 0.58 & $-$0.29 & 0.59 & 0.28 \\
Seasonal & 1.50 & $+$0.60 & 0.61 & 0.42 \\
Irregular & 1.25 & $-$0.14 & 0.63 & 0.37 \\
\bottomrule
\end{tabular}
\caption{Control signal values before and after adaptation (Case 2: Seasonal Amplification).}
\label{tab:case2_policy}
\end{table}


\subsubsection{Design Rationale}

These cases illustrate several key properties of CTRL's test-time adaptation:

\textbf{Component-wise independence.} Each STL component adapts independently, preventing a large trend shift from disrupting well-calibrated seasonal correction. In Case 1, the extreme trend reversal triggers aggressive trend dampening while seasonal adjustment remains moderate.

\textbf{Context-sensitive responses.} The same adaptation mechanism produces opposite adjustments depending on shift direction. Case 1 triggers conservative trend scaling (0.50) with high gate delegation (0.69), while Case 2 triggers aggressive seasonal amplification (1.50) with moderate gate (0.61).

\textbf{Uncertainty-aware delegation.} Confidence parameters absorb uncertainty: high shift scores drop confidence substantially (0.60$\rightarrow$0.21 in Case 1 trend), signaling to the decoder that the LLM is less certain. Simultaneously, gates increase to delegate more correction authority to the learned decoder, implementing a soft handoff under distribution shift.

\textbf{Label-free operation.} Adaptation uses only input statistics from STL decomposition of the lookback window, requiring no ground truth labels. This enables deployment in streaming settings where labels arrive delayed or never.


\subsubsection{Irregular Agent Output}

Unlike trend and seasonal agents that output numerical control signals, the irregular agent generates natural language encoded via frozen GPT-2 and projected to a 32D embedding. Figure~\ref{fig:irregular_output} shows the output format from ETTh2.

\begin{figure}[h!]
\centering
\fbox{\parbox{0.95\columnwidth}{
\textbf{Irregular Agent Output} \\[0.5em]
\textit{[CoT Reasoning]} \textbf{Pattern Recognition:} Historical data shows a slight downward trend with high seasonality. The values fluctuate but generally decrease over time, indicating cyclical behavior. \textbf{Backbone Assessment:} Backbone predictions are slightly off as they predict a stronger downward trend than historical data suggests. They under-predict the mean and volatility compared to historical patterns. \textbf{Uncertainty Consideration:} Confidence should be moderate due to high residual noise (0.2549) indicating significant unexplained variation. The heavy tail difference in kurtosis suggests occasional extreme values. \textbf{Held-out Training Context:} Mean error of 0.0510 and standard deviation of 0.4510, suggesting the backbone under-predicts slightly. \\[0.5em]
\textit{[Reflection]} Given the historical data's slight downward trend and high seasonality, we should DAMPEN backbone predictions (scale $<$ 1.0). A POSITIVE bias is needed since the model under-predicts slightly. A MODERATE gate ($\approx$ 0.5) balances adjustment and uncertainty. \\[0.5em]
\textit{[Control Signal]} Dampen slightly, add positive bias, apply moderate correction strength.
}}
\caption{Chain-of-thought output from irregular agent on ETTh2.}
\label{fig:irregular_output}
\end{figure}

We use text for the irregular component because, unlike trend (characterized by slope) or seasonality (period and amplitude), irregularity lacks a compact parametric description. Phrases like ``high residual noise (0.2549)'' and ``heavy tail difference in kurtosis'' convey nuance that a 4D vector cannot capture. The embedding allows the decoder to learn associations between linguistic descriptions and appropriate corrections.


\section{Control Signal Sensitivity Analysis}
\label{app:sensitivity}

To quantify the decoder's dependence on each control signal component, we perform a parameter sweep analysis on ETTh2 (DLinear, 384$\to$96). For each parameter, we vary its value across the full operational range while holding all other parameters fixed at their LLM-generated values, and measure the mean absolute change in decoder correction output:
\begin{equation}
    \Delta_c = \max_{v \in \mathcal{V}_c} \frac{1}{N} \sum_{i=1}^{N} \left| g_\phi(\hat{y}_i, z) - g_\phi(\hat{y}_i, z_{c \leftarrow v}) \right|
\end{equation}
where $z_{c \leftarrow v}$ denotes the control signal with component $c$ set to sweep value $v$, and $\mathcal{V}_c$ is the operational range for that component. For numeric parameters, we sweep: scale $\in [0.5, 1.5]$, bias $\in [-1.0, 1.0]$, gate $\in [0.0, 1.0]$, confidence $\in [0.0, 1.0]$. For the irregular text embedding, we measure the output change when replacing the LLM-generated text with an empty string, removing the 32-dimensional projected embedding entirely. We use $N = 100$ uniformly sampled test windows.

\begin{table}[H]
\centering
\small
\begin{tabular}{lc}
\toprule
\textbf{Component} & $\boldsymbol{\Delta}$ \textbf{(ETTh2)} \\
\midrule
\textbf{Irregular text emb.} & \textbf{0.032} \\
Trend bias & 0.009 \\
Trend confidence & 0.006 \\
Trend scale & 0.005 \\
Trend gate & 0.005 \\
Seasonal confidence & 0.001 \\
Seasonal scale & 0.001 \\
Seasonal gate & 0.001 \\
Seasonal bias & 0.000 \\
\bottomrule
\end{tabular}
\caption{Control signal sensitivity on ETTh2 (384$\to$96, DLinear). The irregular text embedding is 3.4$\times$ more influential than the most responsive numeric parameter (trend bias).}
\label{tab:sensitivity}
\end{table}

Table~\ref{tab:sensitivity} confirms that the decoder has learned to leverage the 32-dimensional GPT-2 projected text representation as its primary source of semantic context. We additionally test how different text inputs affect the decoder output (Table~\ref{tab:text_variants}), replacing the LLM-generated text with various alternatives.

\begin{table}[H]
\centering
\small
\begin{tabular}{lc}
\toprule
\textbf{Text Variant} & $\boldsymbol{\Delta}$ \\
\midrule
No text (zeroed embedding) & 0.0323 \\
Random unrelated words & 0.0074 \\
Technical jargon & 0.0045 \\
Crisis/volatility description & 0.0010 \\
Random characters & 0.0005 \\
Calm/stable description & 0.0004 \\
\bottomrule
\end{tabular}
\caption{Decoder output change when replacing LLM-generated irregular text with alternative inputs on ETTh2.}
\label{tab:text_variants}
\end{table}

The large gap between removing the text entirely ($\Delta = 0.032$) and substituting alternative texts ($\Delta \leq 0.007$) indicates that the decoder distinguishes the \emph{presence} of the LLM text embedding from its \emph{absence}, while the variation among different text inputs confirms sensitivity to \emph{content}. Combined with the random control signal ablation (Table~\ref{tab:ablation}), which shows random text through the same GPT-2 pipeline performs at the zero-signal level, these results confirm that the decoder utilizes semantic content from the LLM rather than treating the embedding as a generic bias term. This text variant analysis serves as a functional probe of the irregular embedding space: 
by systematically substituting inputs and measuring output divergence, we directly test 
whether the decoder has learned to extract semantic content from the embedding rather 
than relying on its statistical properties alone.


\section{Prompts}
\label{sec:appendix_prompts}

\subsection{Multi-Agent Control Signal Generation}

CTRL employs three specialized agents for STL component-aware control signal generation. Each agent receives multi-scale context (global statistics, few-shot examples, current STL comparison) and outputs structured control signals.

\begin{table}[h!]
\centering
\scriptsize
\setlength{\tabcolsep}{3pt}
\begin{tabular}{p{1.8cm}p{5.2cm}}
\toprule
\multicolumn{2}{l}{\textbf{Trend Specialist Agent}} \\
\midrule
Role & First specialist handling TREND correction \\
Focus & Low-frequency patterns: direction, level shifts, drift, slope accuracy \\
Key Questions & (1) Slope sign match? (2) Slope value match? (3) Systematic over/under-prediction? \\
\midrule
\multicolumn{2}{l}{\textit{Output: 4D control signal}} \\
\texttt{scale} & Trend amplitude adjustment ($>$1 if underestimate, $<$1 if overestimate) \\
\texttt{bias} & Level correction ($+$ if too low, $-$ if too high) \\
\texttt{gate} & Correction aggressiveness \\
\texttt{confidence} & Agent certainty \\
\midrule
\multicolumn{2}{l}{\textit{Dynamic context provided}} \\
Comparison & Historical vs backbone: direction, slope, strength \\
\bottomrule
\end{tabular}
\caption{Trend Agent prompt structure.}
\label{tab:prompt_trend}
\end{table}

\begin{table}[h!]
\centering
\scriptsize
\setlength{\tabcolsep}{3pt}
\begin{tabular}{p{1.8cm}p{5.2cm}}
\toprule
\multicolumn{2}{l}{\textbf{Seasonal Specialist Agent}} \\
\midrule
Role & Second specialist handling SEASONAL correction \\
Focus & High-frequency periodic patterns: daily/weekly cycles, amplitude, phase \\
Key Questions & (1) Correct periodicity? (2) Amplitude correct? (3) Phase aligned? \\
\midrule
\multicolumn{2}{l}{\textit{Output: 4D control signal}} \\
\texttt{scale} & Amplitude adjustment ($>$1 if dampened, $<$1 if exaggerated) \\
\texttt{bias} & Baseline shift (usually $\approx 0$) \\
\texttt{gate} & Correction aggressiveness \\
\texttt{confidence} & Agent certainty \\
\midrule
\multicolumn{2}{l}{\textit{Dynamic context provided}} \\
Comparison & Historical vs backbone: amplitude, dominance, period \\
\bottomrule
\end{tabular}
\caption{Seasonal Agent prompt structure.}
\label{tab:prompt_seasonal}
\end{table}

\begin{table}[h!]
\centering
\scriptsize
\setlength{\tabcolsep}{3pt}
\begin{tabular}{p{1.8cm}p{5.2cm}}
\toprule
\multicolumn{2}{l}{\textbf{Irregular Pattern Agent}} \\
\midrule
Role & Third specialist handling noise/anomalies not explained by trend or seasonality \\
Task & Generate concise text description converted to numerical embedding \\
\midrule
\multicolumn{2}{l}{\textit{Analysis priorities (ordered)}} \\
1. Noise & Variance level: high/low/changing \\
2. Anomalies & Outliers, spikes present? \\
3. Autocorr. & Correlated or white noise? \\
4. Difficulty & Correction predictability \\
5. Distribution & Skewness, tail behavior \\
\midrule
Output & JSON: \texttt{\{"irregular\_analysis": "..."\}} \\
Encoding & Text $\rightarrow$ GPT-2 $\rightarrow$ $d_I$-dim embedding \\
\bottomrule
\end{tabular}
\caption{Irregular Agent prompt structure.}
\label{tab:prompt_irregular}
\end{table}

\subsection{Test-Time Adaptation}

Test-time adaptation detects distribution shift between held-out training and test data using STL decomposition statistics. When shift exceeds threshold $\tau$, agents generate adaptation parameters.

\begin{table}[h!]
\centering
\scriptsize
\setlength{\tabcolsep}{3pt}
\begin{tabular}{p{2.2cm}p{4.8cm}}
\toprule
\multicolumn{2}{l}{\textbf{Trend/Seasonal Adaptation}} \\
\midrule
Input & Component shift analysis: validation slope/dominance vs test, shift ratio (\%), direction change \\
\midrule
\multicolumn{2}{l}{\textit{Output (JSON)}} \\
\texttt{shift\_diagnosis} & Description of detected shift \\
\texttt{impact\_assessment} & Expected effect on predictions \\
\texttt{scale\_multiplier} & Multiplicative adjustment to scale \\
\texttt{bias\_shift} & Additive adjustment to bias \\
\texttt{gate\_multiplier} & Multiplicative adjustment to gate \\
\texttt{confidence\_penalty} & Reduction in confidence $\in [0, 0.5]$ \\
\texttt{reasoning} & Explanation of adaptation strategy \\
\bottomrule
\end{tabular}
\caption{Trend/Seasonal adaptation prompt structure.}
\label{tab:prompt_adapt_ts}
\end{table}

\begin{table}[h!]
\centering
\scriptsize
\setlength{\tabcolsep}{3pt}
\begin{tabular}{p{2.2cm}p{4.8cm}}
\toprule
\multicolumn{2}{l}{\textbf{Irregular Adaptation}} \\
\midrule
Input & Original analysis text, statistics comparison (std, mean, spike ratio, autocorr) for validation vs test \\
\midrule
\multicolumn{2}{l}{\textit{Output (JSON)}} \\
\texttt{needs\_adjustment} & Boolean flag \\
\texttt{updated\_text} & Revised irregular analysis \\
\texttt{adapted\_params} & \{scale, bias, gate, confidence\} adjustments \\
\texttt{reasoning} & Explanation \\
\bottomrule
\end{tabular}
\caption{Irregular adaptation prompt structure.}
\label{tab:prompt_adapt_irreg}
\end{table}

\begin{table}[h!]
\centering
\scriptsize
\setlength{\tabcolsep}{3pt}
\begin{tabular}{p{2.2cm}p{4.8cm}}
\toprule
\multicolumn{2}{l}{\textbf{Unified Distribution Shift Adaptation}} \\
\midrule
\multicolumn{2}{l}{\textit{Input: STL decomposition comparison}} \\
Trend shift & Validation vs test slope, shift ratio, significance \\
Seasonal shift & Validation vs test dominance, shift ratio, significance \\
Raw statistics & Mean, std comparison, overall shift score \\
\midrule
\multicolumn{2}{l}{\textit{Output (JSON)}} \\
Base params & shift\_diagnosis, impact\_assessment, scale\_multiplier, bias\_shift, gate\_multiplier, confidence\_penalty \\
Component adj. & trend\_scale\_adjust, seasonal\_scale\_adjust \\
\texttt{reasoning} & Adaptation rationale \\
\midrule
\multicolumn{2}{l}{\textit{Guidance}} \\
Trend shift & Use bias\_shift for systematic change; trend\_scale\_adjust for magnitude \\
Seasonal shift & Use seasonal\_scale\_adjust for magnitude \\
\bottomrule
\end{tabular}
\caption{Unified adaptation prompt structure.}
\label{tab:prompt_adapt_unified}
\end{table}

